% Diagrama de flujo de información (secuencia) del MVP de ImprontIA (T25, rúbrica 3.2).
% Estudiante -> captura -> backend -> métricas -> indicadores -> docente, con los
% endpoints y tablas reales de TESIS/pfi/mvp al 22/10/2026 (Incremento 16: modelo al cierre)
% (routers/sessions.py::_run_processing lee y escribe la base; processing/ son funciones puras;
% routers/exams.py, models.py, dashboard/src/syncQueue.ts).
% ARCHIVO GENERADO por gen_flujo_informacion.py (misma carpeta): editar allí y regenerar con
%   python3 gen_flujo_informacion.py
% El diagrama se divide en dos partes (fases 1-3 y 4-6) para imprimir el texto a 8 pt
% como mínimo; cada parte se incluye sin escalar en su propia figura:
%   \begin{figure}[p]
%     \centering
%     % Diagrama de flujo de información (secuencia) del MVP de ImprontIA (T25, rúbrica 3.2).
% Estudiante -> captura -> backend -> métricas -> indicadores -> docente, con los
% endpoints y tablas reales de TESIS/pfi/mvp al 22/10/2026 (Incremento 16: modelo al cierre)
% (routers/sessions.py::_run_processing lee y escribe la base; processing/ son funciones puras;
% routers/exams.py, models.py, dashboard/src/syncQueue.ts).
% ARCHIVO GENERADO por gen_flujo_informacion.py (misma carpeta): editar allí y regenerar con
%   python3 gen_flujo_informacion.py
% El diagrama se divide en dos partes (fases 1-3 y 4-6) para imprimir el texto a 8 pt
% como mínimo; cada parte se incluye sin escalar en su propia figura:
%   \begin{figure}[p]
%     \centering
%     % Diagrama de flujo de información (secuencia) del MVP de ImprontIA (T25, rúbrica 3.2).
% Estudiante -> captura -> backend -> métricas -> indicadores -> docente, con los
% endpoints y tablas reales de TESIS/pfi/mvp al 22/10/2026 (Incremento 16: modelo al cierre)
% (routers/sessions.py::_run_processing lee y escribe la base; processing/ son funciones puras;
% routers/exams.py, models.py, dashboard/src/syncQueue.ts).
% ARCHIVO GENERADO por gen_flujo_informacion.py (misma carpeta): editar allí y regenerar con
%   python3 gen_flujo_informacion.py
% El diagrama se divide en dos partes (fases 1-3 y 4-6) para imprimir el texto a 8 pt
% como mínimo; cada parte se incluye sin escalar en su propia figura:
%   \begin{figure}[p]
%     \centering
%     % Diagrama de flujo de información (secuencia) del MVP de ImprontIA (T25, rúbrica 3.2).
% Estudiante -> captura -> backend -> métricas -> indicadores -> docente, con los
% endpoints y tablas reales de TESIS/pfi/mvp al 22/10/2026 (Incremento 16: modelo al cierre)
% (routers/sessions.py::_run_processing lee y escribe la base; processing/ son funciones puras;
% routers/exams.py, models.py, dashboard/src/syncQueue.ts).
% ARCHIVO GENERADO por gen_flujo_informacion.py (misma carpeta): editar allí y regenerar con
%   python3 gen_flujo_informacion.py
% El diagrama se divide en dos partes (fases 1-3 y 4-6) para imprimir el texto a 8 pt
% como mínimo; cada parte se incluye sin escalar en su propia figura:
%   \begin{figure}[p]
%     \centering
%     \input{figures/diagramas/flujo_informacion}   % o flujo_informacion_b
%     \caption{...}\label{fig:diag_flujo}
%   \end{figure}
% IMPORTANTE: las bibliotecas TikZ se cargan UNA vez en el preámbulo de main_full.tex:
%   \usetikzlibrary{arrows.meta,positioning,fit,calc,backgrounds,shapes.multipart}
% (cargarlas dentro de un entorno figure no sirve: la carga queda local a ese grupo
% y el segundo diagrama falla con "Unknown arrow tip kind 'Stealth'").
\begin{tikzpicture}[
    font={\linespread{1}\scriptsize},
    execute at begin node={\hyphenpenalty=10000},
    actor/.style={draw=accent, fill=accentlight, rounded corners=1.5pt, text width=22mm,
                  align=center, minimum height=9mm, inner sep=2pt, font={\linespread{1}\scriptsize\bfseries}},
    vida/.style={draw=black!35, dashed},
    msg/.style={-{Stealth[length=1.6mm]}, draw=black!80, semithick},
    ret/.style={-{Stealth[length=1.6mm]}, draw=black!60, dashed},
    lbl/.style={font={\linespread{1}\scriptsize}, fill=white, inner sep=1pt, align=center},
    nota/.style={draw=black!45, fill=yellow!8, font={\linespread{1}\scriptsize}, align=left, inner sep=2pt, text width=#1},
    nota/.default=25mm,
    fase/.style={font={\linespread{1}\scriptsize\bfseries}, text=white, fill=accent, inner sep=2pt, anchor=north west},
    bucle/.style={font={\linespread{1}\scriptsize\bfseries}, draw=black!55, fill=white, anchor=north west, inner sep=1.5pt},
]
% ---- Participantes y líneas de vida ----

\node[actor] at (0.0,0) {Estudiante};
\node[actor] at (2.65,0) {Navegador\\ \texttt{TakeExam.tsx}\\ + \texttt{capture.ts}};
\node[actor] at (5.3,0) {API FastAPI\\ \texttt{routers/}};
\node[actor] at (7.95,0) {Procesamiento\\ \texttt{processing/}\\ (funciones puras)};
\node[actor] at (10.6,0) {Base de datos};
\node[actor] at (13.25,0) {Docente\\ (dashboard)};
\foreach \x in {0.0,2.65,5.3,7.95,10.6,13.25} { \draw[vida] (\x,-0.55) -- (\x,-12.72); }
% ---- Mensajes (generados; ver comentario de cabecera) ----
\draw[accent, thin] (-1.2,-0.87) -- (14.45,-0.87);
\node[fase] at (-1.2,-0.87) {1. Información previa al consentimiento (RF01, RF02)};
\draw[msg] (0.0,-1.98) -- node[lbl, above] {abre \texttt{/exam/\{código\}}} (2.65,-1.98);
\draw[msg] (2.65,-2.63) -- node[lbl, above] {\texttt{GET /exams/\{código\}/info}} (5.3,-2.63);
\draw[msg] (5.3,-3.28) -- node[lbl, above] {lee la política vigente} (10.6,-3.28);
\draw[ret] (5.3,-3.93) -- node[lbl, above] {qué se registra + señales + reglas (sin consigna)} (2.65,-3.93);
\draw[msg] (0.0,-4.58) -- node[lbl, above] {identificador + consentimiento} (2.65,-4.58);
\draw[accent, thin] (-1.2,-4.94) -- (14.45,-4.94);
\node[fase] at (-1.2,-4.94) {2. Inicio de la sesión (RF03, RNF04)};
\draw[msg] (2.65,-6.05) -- node[lbl, above] {\texttt{POST /sessions/start}} (5.3,-6.05);
\draw[msg] (5.3,-6.70) -- node[lbl, above] {\texttt{Session}: fija la versión de política} (10.6,-6.70);
\draw[ret] (5.3,-7.35) -- node[lbl, above] {token de sesión, señales habilitadas, consigna} (2.65,-7.35);
\draw[accent, thin] (-1.2,-7.71) -- (14.45,-7.71);
\node[fase] at (-1.2,-7.71) {3. Captura durante la sesión (RF04, RNF01, RNF07)};
\node[bucle] at (-0.95,-8.28) {bucle [cada 5\,s; ante fallas de red, con \textit{backoff}]};
\draw[msg] (0.0,-9.31) -- node[lbl, above] {teclado, mouse, foco, pegado} (2.65,-9.31);
\node[nota=35mm, anchor=north west] at (2.77,-9.49) {filtra por señales habilitadas; conserva solo tiempos y conteos, nunca el texto; buffer local acotado};
\draw[msg] (2.65,-11.59) -- node[lbl, above] {\texttt{POST /sessions/\{id\}/events}\\ (\texttt{seq} del lote + eventos)} (5.3,-11.59);
\draw[msg] (5.3,-12.24) -- node[lbl, above] {\texttt{Event} (solo inserción; un lote repetido se ignora)} (10.6,-12.24);
\draw[black!55, thin] (-0.95,-8.28) rectangle (11.15,-12.42);
\end{tikzpicture}
   % o flujo_informacion_b
%     \caption{...}\label{fig:diag_flujo}
%   \end{figure}
% IMPORTANTE: las bibliotecas TikZ se cargan UNA vez en el preámbulo de main_full.tex:
%   \usetikzlibrary{arrows.meta,positioning,fit,calc,backgrounds,shapes.multipart}
% (cargarlas dentro de un entorno figure no sirve: la carga queda local a ese grupo
% y el segundo diagrama falla con "Unknown arrow tip kind 'Stealth'").
\begin{tikzpicture}[
    font={\linespread{1}\scriptsize},
    execute at begin node={\hyphenpenalty=10000},
    actor/.style={draw=accent, fill=accentlight, rounded corners=1.5pt, text width=22mm,
                  align=center, minimum height=9mm, inner sep=2pt, font={\linespread{1}\scriptsize\bfseries}},
    vida/.style={draw=black!35, dashed},
    msg/.style={-{Stealth[length=1.6mm]}, draw=black!80, semithick},
    ret/.style={-{Stealth[length=1.6mm]}, draw=black!60, dashed},
    lbl/.style={font={\linespread{1}\scriptsize}, fill=white, inner sep=1pt, align=center},
    nota/.style={draw=black!45, fill=yellow!8, font={\linespread{1}\scriptsize}, align=left, inner sep=2pt, text width=#1},
    nota/.default=25mm,
    fase/.style={font={\linespread{1}\scriptsize\bfseries}, text=white, fill=accent, inner sep=2pt, anchor=north west},
    bucle/.style={font={\linespread{1}\scriptsize\bfseries}, draw=black!55, fill=white, anchor=north west, inner sep=1.5pt},
]
% ---- Participantes y líneas de vida ----

\node[actor] at (0.0,0) {Estudiante};
\node[actor] at (2.65,0) {Navegador\\ \texttt{TakeExam.tsx}\\ + \texttt{capture.ts}};
\node[actor] at (5.3,0) {API FastAPI\\ \texttt{routers/}};
\node[actor] at (7.95,0) {Procesamiento\\ \texttt{processing/}\\ (funciones puras)};
\node[actor] at (10.6,0) {Base de datos};
\node[actor] at (13.25,0) {Docente\\ (dashboard)};
\foreach \x in {0.0,2.65,5.3,7.95,10.6,13.25} { \draw[vida] (\x,-0.55) -- (\x,-12.72); }
% ---- Mensajes (generados; ver comentario de cabecera) ----
\draw[accent, thin] (-1.2,-0.87) -- (14.45,-0.87);
\node[fase] at (-1.2,-0.87) {1. Información previa al consentimiento (RF01, RF02)};
\draw[msg] (0.0,-1.98) -- node[lbl, above] {abre \texttt{/exam/\{código\}}} (2.65,-1.98);
\draw[msg] (2.65,-2.63) -- node[lbl, above] {\texttt{GET /exams/\{código\}/info}} (5.3,-2.63);
\draw[msg] (5.3,-3.28) -- node[lbl, above] {lee la política vigente} (10.6,-3.28);
\draw[ret] (5.3,-3.93) -- node[lbl, above] {qué se registra + señales + reglas (sin consigna)} (2.65,-3.93);
\draw[msg] (0.0,-4.58) -- node[lbl, above] {identificador + consentimiento} (2.65,-4.58);
\draw[accent, thin] (-1.2,-4.94) -- (14.45,-4.94);
\node[fase] at (-1.2,-4.94) {2. Inicio de la sesión (RF03, RNF04)};
\draw[msg] (2.65,-6.05) -- node[lbl, above] {\texttt{POST /sessions/start}} (5.3,-6.05);
\draw[msg] (5.3,-6.70) -- node[lbl, above] {\texttt{Session}: fija la versión de política} (10.6,-6.70);
\draw[ret] (5.3,-7.35) -- node[lbl, above] {token de sesión, señales habilitadas, consigna} (2.65,-7.35);
\draw[accent, thin] (-1.2,-7.71) -- (14.45,-7.71);
\node[fase] at (-1.2,-7.71) {3. Captura durante la sesión (RF04, RNF01, RNF07)};
\node[bucle] at (-0.95,-8.28) {bucle [cada 5\,s; ante fallas de red, con \textit{backoff}]};
\draw[msg] (0.0,-9.31) -- node[lbl, above] {teclado, mouse, foco, pegado} (2.65,-9.31);
\node[nota=35mm, anchor=north west] at (2.77,-9.49) {filtra por señales habilitadas; conserva solo tiempos y conteos, nunca el texto; buffer local acotado};
\draw[msg] (2.65,-11.59) -- node[lbl, above] {\texttt{POST /sessions/\{id\}/events}\\ (\texttt{seq} del lote + eventos)} (5.3,-11.59);
\draw[msg] (5.3,-12.24) -- node[lbl, above] {\texttt{Event} (solo inserción; un lote repetido se ignora)} (10.6,-12.24);
\draw[black!55, thin] (-0.95,-8.28) rectangle (11.15,-12.42);
\end{tikzpicture}
   % o flujo_informacion_b
%     \caption{...}\label{fig:diag_flujo}
%   \end{figure}
% IMPORTANTE: las bibliotecas TikZ se cargan UNA vez en el preámbulo de main_full.tex:
%   \usetikzlibrary{arrows.meta,positioning,fit,calc,backgrounds,shapes.multipart}
% (cargarlas dentro de un entorno figure no sirve: la carga queda local a ese grupo
% y el segundo diagrama falla con "Unknown arrow tip kind 'Stealth'").
\begin{tikzpicture}[
    font={\linespread{1}\scriptsize},
    execute at begin node={\hyphenpenalty=10000},
    actor/.style={draw=accent, fill=accentlight, rounded corners=1.5pt, text width=22mm,
                  align=center, minimum height=9mm, inner sep=2pt, font={\linespread{1}\scriptsize\bfseries}},
    vida/.style={draw=black!35, dashed},
    msg/.style={-{Stealth[length=1.6mm]}, draw=black!80, semithick},
    ret/.style={-{Stealth[length=1.6mm]}, draw=black!60, dashed},
    lbl/.style={font={\linespread{1}\scriptsize}, fill=white, inner sep=1pt, align=center},
    nota/.style={draw=black!45, fill=yellow!8, font={\linespread{1}\scriptsize}, align=left, inner sep=2pt, text width=#1},
    nota/.default=25mm,
    fase/.style={font={\linespread{1}\scriptsize\bfseries}, text=white, fill=accent, inner sep=2pt, anchor=north west},
    bucle/.style={font={\linespread{1}\scriptsize\bfseries}, draw=black!55, fill=white, anchor=north west, inner sep=1.5pt},
]
% ---- Participantes y líneas de vida ----

\node[actor] at (0.0,0) {Estudiante};
\node[actor] at (2.65,0) {Navegador\\ \texttt{TakeExam.tsx}\\ + \texttt{capture.ts}};
\node[actor] at (5.3,0) {API FastAPI\\ \texttt{routers/}};
\node[actor] at (7.95,0) {Procesamiento\\ \texttt{processing/}\\ (funciones puras)};
\node[actor] at (10.6,0) {Base de datos};
\node[actor] at (13.25,0) {Docente\\ (dashboard)};
\foreach \x in {0.0,2.65,5.3,7.95,10.6,13.25} { \draw[vida] (\x,-0.55) -- (\x,-12.72); }
% ---- Mensajes (generados; ver comentario de cabecera) ----
\draw[accent, thin] (-1.2,-0.87) -- (14.45,-0.87);
\node[fase] at (-1.2,-0.87) {1. Información previa al consentimiento (RF01, RF02)};
\draw[msg] (0.0,-1.98) -- node[lbl, above] {abre \texttt{/exam/\{código\}}} (2.65,-1.98);
\draw[msg] (2.65,-2.63) -- node[lbl, above] {\texttt{GET /exams/\{código\}/info}} (5.3,-2.63);
\draw[msg] (5.3,-3.28) -- node[lbl, above] {lee la política vigente} (10.6,-3.28);
\draw[ret] (5.3,-3.93) -- node[lbl, above] {qué se registra + señales + reglas (sin consigna)} (2.65,-3.93);
\draw[msg] (0.0,-4.58) -- node[lbl, above] {identificador + consentimiento} (2.65,-4.58);
\draw[accent, thin] (-1.2,-4.94) -- (14.45,-4.94);
\node[fase] at (-1.2,-4.94) {2. Inicio de la sesión (RF03, RNF04)};
\draw[msg] (2.65,-6.05) -- node[lbl, above] {\texttt{POST /sessions/start}} (5.3,-6.05);
\draw[msg] (5.3,-6.70) -- node[lbl, above] {\texttt{Session}: fija la versión de política} (10.6,-6.70);
\draw[ret] (5.3,-7.35) -- node[lbl, above] {token de sesión, señales habilitadas, consigna} (2.65,-7.35);
\draw[accent, thin] (-1.2,-7.71) -- (14.45,-7.71);
\node[fase] at (-1.2,-7.71) {3. Captura durante la sesión (RF04, RNF01, RNF07)};
\node[bucle] at (-0.95,-8.28) {bucle [cada 5\,s; ante fallas de red, con \textit{backoff}]};
\draw[msg] (0.0,-9.31) -- node[lbl, above] {teclado, mouse, foco, pegado} (2.65,-9.31);
\node[nota=35mm, anchor=north west] at (2.77,-9.49) {filtra por señales habilitadas; conserva solo tiempos y conteos, nunca el texto; buffer local acotado};
\draw[msg] (2.65,-11.59) -- node[lbl, above] {\texttt{POST /sessions/\{id\}/events}\\ (\texttt{seq} del lote + eventos)} (5.3,-11.59);
\draw[msg] (5.3,-12.24) -- node[lbl, above] {\texttt{Event} (solo inserción; un lote repetido se ignora)} (10.6,-12.24);
\draw[black!55, thin] (-0.95,-8.28) rectangle (11.15,-12.42);
\end{tikzpicture}
   % o flujo_informacion_b
%     \caption{...}\label{fig:diag_flujo}
%   \end{figure}
% IMPORTANTE: las bibliotecas TikZ se cargan UNA vez en el preámbulo de main_full.tex:
%   \usetikzlibrary{arrows.meta,positioning,fit,calc,backgrounds,shapes.multipart}
% (cargarlas dentro de un entorno figure no sirve: la carga queda local a ese grupo
% y el segundo diagrama falla con "Unknown arrow tip kind 'Stealth'").
\begin{tikzpicture}[
    font={\linespread{1}\scriptsize},
    execute at begin node={\hyphenpenalty=10000},
    actor/.style={draw=accent, fill=accentlight, rounded corners=1.5pt, text width=22mm,
                  align=center, minimum height=9mm, inner sep=2pt, font={\linespread{1}\scriptsize\bfseries}},
    vida/.style={draw=black!35, dashed},
    msg/.style={-{Stealth[length=1.6mm]}, draw=black!80, semithick},
    ret/.style={-{Stealth[length=1.6mm]}, draw=black!60, dashed},
    lbl/.style={font={\linespread{1}\scriptsize}, fill=white, inner sep=1pt, align=center},
    nota/.style={draw=black!45, fill=yellow!8, font={\linespread{1}\scriptsize}, align=left, inner sep=2pt, text width=#1},
    nota/.default=25mm,
    fase/.style={font={\linespread{1}\scriptsize\bfseries}, text=white, fill=accent, inner sep=2pt, anchor=north west},
    bucle/.style={font={\linespread{1}\scriptsize\bfseries}, draw=black!55, fill=white, anchor=north west, inner sep=1.5pt},
]
% ---- Participantes y líneas de vida ----

\node[actor] at (0.0,0) {Estudiante};
\node[actor] at (2.65,0) {Navegador\\ \texttt{TakeExam.tsx}\\ + \texttt{capture.ts}};
\node[actor] at (5.3,0) {API FastAPI\\ \texttt{routers/}};
\node[actor] at (7.95,0) {Procesamiento\\ \texttt{processing/}\\ (funciones puras)};
\node[actor] at (10.6,0) {Base de datos};
\node[actor] at (13.25,0) {Docente\\ (dashboard)};
\foreach \x in {0.0,2.65,5.3,7.95,10.6,13.25} { \draw[vida] (\x,-0.55) -- (\x,-12.72); }
% ---- Mensajes (generados; ver comentario de cabecera) ----
\draw[accent, thin] (-1.2,-0.87) -- (14.45,-0.87);
\node[fase] at (-1.2,-0.87) {1. Información previa al consentimiento (RF01, RF02)};
\draw[msg] (0.0,-1.98) -- node[lbl, above] {abre \texttt{/exam/\{código\}}} (2.65,-1.98);
\draw[msg] (2.65,-2.63) -- node[lbl, above] {\texttt{GET /exams/\{código\}/info}} (5.3,-2.63);
\draw[msg] (5.3,-3.28) -- node[lbl, above] {lee la política vigente} (10.6,-3.28);
\draw[ret] (5.3,-3.93) -- node[lbl, above] {qué se registra + señales + reglas (sin consigna)} (2.65,-3.93);
\draw[msg] (0.0,-4.58) -- node[lbl, above] {identificador + consentimiento} (2.65,-4.58);
\draw[accent, thin] (-1.2,-4.94) -- (14.45,-4.94);
\node[fase] at (-1.2,-4.94) {2. Inicio de la sesión (RF03, RNF04)};
\draw[msg] (2.65,-6.05) -- node[lbl, above] {\texttt{POST /sessions/start}} (5.3,-6.05);
\draw[msg] (5.3,-6.70) -- node[lbl, above] {\texttt{Session}: fija la versión de política} (10.6,-6.70);
\draw[ret] (5.3,-7.35) -- node[lbl, above] {token de sesión, señales habilitadas, consigna} (2.65,-7.35);
\draw[accent, thin] (-1.2,-7.71) -- (14.45,-7.71);
\node[fase] at (-1.2,-7.71) {3. Captura durante la sesión (RF04, RNF01, RNF07)};
\node[bucle] at (-0.95,-8.28) {bucle [cada 5\,s; ante fallas de red, con \textit{backoff}]};
\draw[msg] (0.0,-9.31) -- node[lbl, above] {teclado, mouse, foco, pegado} (2.65,-9.31);
\node[nota=35mm, anchor=north west] at (2.77,-9.49) {filtra por señales habilitadas; conserva solo tiempos y conteos, nunca el texto; buffer local acotado};
\draw[msg] (2.65,-11.59) -- node[lbl, above] {\texttt{POST /sessions/\{id\}/events}\\ (\texttt{seq} del lote + eventos)} (5.3,-11.59);
\draw[msg] (5.3,-12.24) -- node[lbl, above] {\texttt{Event} (solo inserción; un lote repetido se ignora)} (10.6,-12.24);
\draw[black!55, thin] (-0.95,-8.28) rectangle (11.15,-12.42);
\end{tikzpicture}
